% TOP_PROBLEM: {"id": 30004922, "title": "Derandomizing Space-Bounded Computation", "category_id": 15, "category": {"id": 15, "name": "computer_science", "display_name": "Computer Science", "description": "Computational complexity, algorithms, and theoretical CS.", "slug": "computer-science", "order_index": 15, "created_at": "2026-07-31T15:26:25.671Z"}, "status": "open"}
% FIRST_POSTED: 2026-09-27
% !TeX program = lualatex
% Compile twice: lualatex open_problems_500.tex
% All references and prose are embedded; no BibTeX database or external inputs.
\documentclass[11pt,a4paper]{article}
\usepackage[margin=24mm,headheight=14pt]{geometry}
\usepackage{fontspec}
\setmainfont{Latin Modern Roman}
\setsansfont{Latin Modern Sans}
\setmonofont{Latin Modern Mono}
\usepackage{amsmath,amssymb,mathtools}
\usepackage{unicode-math}
\setmathfont{Latin Modern Math}
\usepackage{microtype}
\usepackage{enumitem,xurl,needspace}
\usepackage[dvipsnames]{xcolor}
\usepackage{hyperref}
\hypersetup{unicode=true,colorlinks=true,linkcolor=MidnightBlue,urlcolor=MidnightBlue,
 pdftitle={500 Mathematical Problem Statements},pdfauthor={Source-annotated compilation},bookmarksnumbered=true}
\usepackage{bookmark}
\usepackage{tocloft}
\setlength{\cftsecnumwidth}{3em}
\cftsetpnumwidth{2.5em}
\cftsetrmarg{4em}
\usepackage{titlesec}
\titleformat{\section}{\Large\bfseries\raggedright}{\thesection}{0.7em}{}
\usepackage{fancyhdr}
\pagestyle{fancy}\fancyhf{}
\fancyhead[L]{\small\sffamily Mathematical problem statements}
\fancyhead[R]{\small Source edition 22}
\fancyfoot[C]{\thepage}
\setlength{\parindent}{0pt}
\setlength{\parskip}{5pt plus 1pt}
\setlength{\emergencystretch}{3em}
\setcounter{tocdepth}{1}
\setcounter{secnumdepth}{1}
\setlist[itemize]{leftmargin=3em,itemsep=5pt,topsep=4pt,parsep=0pt}
\allowdisplaybreaks
\newcommand{\N}{\mathbb{N}}
\newcommand{\Z}{\mathbb{Z}}
\newcommand{\Q}{\mathbb{Q}}
\newcommand{\R}{\mathbb{R}}
\newcommand{\C}{\mathbb{C}}
\newcommand{\F}{\mathbb{F}}
\newcommand{\A}{\mathbb{A}}
\newcommand{\PP}{\mathbb P}
\newcommand{\E}{\mathbb{E}}
\newcommand{\eps}{\varepsilon}
\newcommand{\dd}{\,\mathrm d}
\newcommand{\sref}[2]{\hyperlink{source:#1:#2}{[S#2]}}
\newcommand{\eref}[2]{\hyperlink{extra:#1:#2}{[E#2]}}
% Turn the next switch off for a shorter reading copy. The quotations preserve
% every exactTarget field, including qualifications omitted from a short summary.
\newif\ifshowcatalogue
\showcataloguetrue
\newcommand{\cataloguescope}[1]{\ifshowcatalogue
 \paragraph{Exact scope in the supplied catalogue.}
 \begin{quote}\small #1\end{quote}\fi}
% Compile twice with: lualatex 30004922.tex
\numberwithin{equation}{section}
\hypersetup{pdftitle={Research attempt -- problem 30004922},pdfauthor={Research notebook; original compilation retained}}
\fancyhead[L]{\small\sffamily Research notebook}
\fancyhead[R]{\small\sffamily Problem 30004922 | 27 September 2026}
\begin{document}
\setcounter{section}{0}
% ================================================================
% problemId: 30004922
% Canonical problem ID: 30004922
% exactTarget SHA-256: 30bae69d2d4994611f4de40617344461d615b0d70538263d7b41fca724c6cf1d
\clearpage
\section*{\texorpdfstring{Derandomizing Space-Bounded Computation}{Derandomizing Space-Bounded Computation}}\label{30004922}
\subsection{Definitions and mathematical statement}
A probabilistic logarithmic-space machine has read-only input, $O(\log n)$ work bits, random bits, and polynomial running time. The class $\mathrm{BPL}$ consists of languages it decides with acceptance probability at least $2/3$ on yes inputs and at most $1/3$ on no inputs. The deterministic class $\mathrm L$ uses the same logarithmic workspace without randomness. Determine whether
\[
 \mathrm{BPL}=\mathrm L.
\]
All bounds are uniform over inputs of each length. Space counts writable working storage, not the read-only input or the number of random bits ever read. The question is exact derandomization of the language class, not merely a small increase in the workspace bound.
\par\smallskip\textit{Formulation and concept references:} \sref{30004922}{1}.
\cataloguescope{Prove or disprove that every language decidable by a probabilistic Turing machine using O(log n) workspace, polynomial time, and bounded two-sided error is decidable by a deterministic Turing machine using O(log n) workspace.}
\subsection{Short English statement}
Can randomness always be removed from a polynomial-time decision procedure using logarithmic memory without increasing its memory requirement beyond a constant factor?
% BEGIN RESEARCH ADDITION ATTEMPT-163
\subsection{Research attempt}

\paragraph{Current frontier.}
The exact target is a uniform $O(\log n)$ workspace bound. Hoza's general
simulation improves the Saks--Zhou bound to
$O(\log^{3/2}n/\sqrt{\log\log n})$, not to logarithmic space
\eref{30004922}{1}. Cheng--Wu's newer weighted generators improve several
parameter regimes, including short-wide and regular branching programs;
those hypotheses do not cover all polynomial-length, polynomial-width
programs arising from BPL \sref{30004922}{1}, Sections 1.1--1.4. Hitting sets
can suffice for two-sided derandomization through a nontrivial reduction,
as Cheng--Hoza prove \eref{30004922}{2}. None of these results licenses
replacing an arbitrary acceptance probability by an easily computed
unweighted average over a logarithmic seed.

\paragraph{Attack route.}
There are three plausible ways to avoid retaining a long random seed:
compute an approximate backward value function, forget old randomness
through contraction, or use a signed short-seed quadrature. The first
requires a uniform way to generate local approximate values. The second
requires quantitative loss of memory in the acceptance observable. The
third requires both short descriptions and control of cancellation.
I pursue a common error-budget formulation, then try to break each
shortcut on elementary examples.

\paragraph{Attempt: an occupancy-weighted residual identity.}
Fix an input of length $n$ and pad the probabilistic machine to exactly
$T\le n^c$ steps. Its configurations fit in $O(\log n)$ bits. By adding
a layer index, it becomes a read-once branching program with stochastic
transition matrices $P_0,\ldots,P_{T-1}$ on at most $n^c$ states; each
row can be computed in logarithmic space. The terminal function $f$ is
the indicator of acceptance. Let $X_i$ be the random state at layer $i$.

For arbitrary real functions $v_i$ on the layers define
\[
 r_i(v)=v_i(v)-(P_iv_{i+1})(v),\qquad 0\le i<T.
\]
The following exact identity is useful:
\begin{equation}\label{30004922:residual}
 v_0(x_0)-\E f(X_T)
 =\sum_{i=0}^{T-1}\E r_i(X_i)
       +\E\bigl(v_T(X_T)-f(X_T)\bigr).
\end{equation}
It follows by summing
$\E v_i(X_i)-\E v_{i+1}(X_{i+1})=\E r_i(X_i)$.
Everything is finite. In particular,
\begin{equation}\label{30004922:residual-bound}
 |v_0(x_0)-\Pr(\mathrm{accept})|
 \le\sum_{i<T}\|r_i\|_\infty+\|v_T-f\|_\infty.
\end{equation}
This distinguishes an exact occupancy-weighted error budget from a
potentially much worse uniform local budget.

\textbf{Sufficient value-generator lemma.} Suppose that, uniformly for
every such program, a deterministic logarithmic-space algorithm outputs
$v_i(v)$ as a dyadic rational with $O(\log n)$ bits, and that
\[
 \|v_T-f\|_\infty\le1/24,
 \qquad\sum_{i<T}\|r_i\|_\infty\le1/24.
\]
Then BPL is contained in L: compute only $v_0(x_0)$ and compare it with
$1/2$. The error at most $1/12$ separates acceptance probabilities at
least $2/3$ from those at most $1/3$. The reverse inclusion is immediate.
No table of all configurations is stored in this algorithm.

The hypothesis is not obtained simply by rounding the true backward
probabilities. Such rounding does show that suitable polynomial-size
tables exist: with denominator $2^b$, $b=O(\log T)$, one can make each
residual $O(1/T)$. It does \emph{not} supply a logspace generator for their
entries. A polynomial-size table of dyadic values can be checked in
logspace when provided on a separate read-only certificate tape, by
looping over its rows. That is a different resource from computing the
table deterministically. Treating the certificate as free advice would
violate the target's uniformity.

\paragraph{Attempt: a genuine contraction subclass.}
Let $\ell\le C\log n$ be an explicitly specified number of final random
steps. Suppose that the laws of the terminal state, started at any two
states $u,v$ of layer $T-\ell$, differ in total variation by at most
$1/12$. Then the acceptance probability is deterministically
approximable within $1/12$ in logarithmic space: select any fixed state
of that layer, enumerate the $2^\ell$ suffixes of random bits, simulate
them, and count accepting suffixes. A suffix, a configuration, the loop
index, and the counter each occupy $O(\log n)$ bits. Time is polynomial.
The true distribution at layer $T-\ell$ is a mixture of the possible
starting states, and the total-variation hypothesis bounds its difference
from this fixed-state computation.

This criterion is slightly stronger than needed: it suffices that the
\emph{acceptance probabilities} from those states differ by at most
$1/12$, rather than their full terminal laws. However, certifying that
weaker condition by first computing all acceptance probabilities is
circular. The criterion is a solved subclass, not an argument that every
branching program has a rapidly contracting suffix.

\paragraph{Attempt: signed quadrature and bit complexity.}
Here is another precise sufficient input. Suppose a uniform logarithmic-
space procedure generates, for $r\in\{0,1\}^d$, a $T$-bit string $G(r)$
one bit at a time and a weight $w_r=z_r2^{-b}$, where
\[
 d,b\le C\log n,\qquad |z_r|\le n^C,
 \qquad
 \left|\Pr(\mathrm{accept})-\sum_r w_r f(G(r))\right|<1/12.
\]
Enumerating all seeds and accumulating the integers $z_rf(G(r))$ uses
only $O(\log n)$ space: the absolute numerator is at most
$2^d n^C$, and the common denominator has $O(\log n)$ bits. The generated
bits can be recomputed whenever needed. This would again yield BPL=L.
The inclusion of the weight encoding is essential; a short seed alone
does not bound the space needed for exponentially precise signed weights.
Weighted generators are an established framework \sref{30004922}{1}; the
specific unrestricted logarithmic parameters above are the missing input,
not an application of a cited theorem.

Can local signed compressions be composed without increasing that cost?
For matrices with the operator norm induced by $\ell^\infty$, telescoping
products gives
\begin{equation}\label{30004922:product-error}
 \|P_0\cdots P_{T-1}-Q_0\cdots Q_{T-1}\|_\infty
 \le\sum_{i=0}^{T-1}
       \|P_i-Q_i\|_\infty
       \prod_{j<i}\|Q_j\|_\infty.
\end{equation}
Indeed, insert $Q_0, Q_1,\ldots$ successively on the left and keep the
remaining $P$ matrices on the right; stochastic matrices have norm one.
Thus a controlled sum of local errors is insufficient when signed
intermediate matrices have large absolute row sums.

\paragraph{Stress test: three explicit failures.}
First, an absorbing accepting state and an absorbing rejecting state
never mix: their terminal laws remain at total-variation distance one.
This can happen in a deterministic logspace computation. Therefore
rapid mixing is not even a necessary structural property of problems in
L, and cannot be asserted for all BPL computations without a further
transformation that preserves the answer.

Second, preservation of constants is not the same as norm stability.
For $a>0$,
\[
 Q=\begin{pmatrix}1+a&-a\\-a&1+a\end{pmatrix}
 \quad\text{satisfies}\quad
 Q\binom11=\binom11,
 \qquad Q^t\binom1{-1}=(1+2a)^t\binom1{-1}.
\]
Every row sums to one, yet errors in the antisymmetric direction are
exponentially amplified. At $a=1/4$, twelve applications amplify that
direction by $(3/2)^{12}$, verified exactly in the script. This is a
counterexample to the proposed \emph{stability lemma}, not to BPL=L.

Third, even a recursion of only $O(\log n)$ levels can multiply
polynomial absolute-weight budgets into $n^{O(\log n)}$. Directly
accumulating such numerators may require $\Theta(\log^2 n)$ bits.
Avoiding this cost would need a compensated composition, a norm in which
the intermediate operators remain controlled, or a different use of
hitting sets such as the framework of \eref{30004922}{2}. I have not derived
such a uniform composition rule with logarithmic workspace.

The value-generator lemma is also adversarially checked for rounding:
per-layer errors of a fixed constant would sum to order $T$, and so
cannot be used. The required $O(1/T)$ scale is cheap to \emph{encode},
but not thereby cheap to \emph{compute}. That is the exact point where
the attempted conversion of local verification into a deterministic
simulation stops.

\paragraph{Outcome.}
\textbf{Outcome: Exploratory approach with unresolved gap.}

The residual identity, the contraction-subclass simulation, the signed-
quadrature bit bound, and the product-error estimate are proved. They
provide compatible tests for a proposed derandomization mechanism and
exhibit concrete failures of memory erasure and signed-weight stability.
No general logspace value generator or stable logarithmic-seed quadrature
has been constructed. The additional input needed for either would
already imply the full language-class equality, not merely a small
improvement in workspace.

% Keep the local source heading and initial references together.
\Needspace{12\baselineskip}
% END RESEARCH ADDITION ATTEMPT-163
\subsection{Sources}
\begin{itemize}\raggedright
\item[\textnormal{[S1]}]\hypertarget{source:30004922:1}{} Kuan Cheng and Ruiyang Wu, Weighted Pseudorandom Generators for Read-Once Branching Programs via Weighted Pseudorandom Reductions, Introduction.
\url{https://arxiv.org/html/2502.08272v5}.
{\small\textit{Catalogue formulation source.}}
% BEGIN RESEARCH ADDITION SOURCES-163
\item[\textnormal{[E1]}]\hypertarget{extra:30004922:1}{} William M. Hoza, ``Better Pseudodistributions and Derandomization for Space-Bounded Computation,'' \emph{APPROX/RANDOM 2021}, LIPIcs 207, article 28, 28:1--28:23. Main generator and deterministic-space consequences.
\url{https://doi.org/10.4230/LIPIcs.APPROX/RANDOM.2021.28}.
{\small\textit{Author summary and proceedings version distinguish the improved bound from logarithmic space.}}
\item[\textnormal{[E2]}]\hypertarget{extra:30004922:2}{} Kuan Cheng and William M. Hoza, ``Hitting Sets Give Two-Sided Derandomization of Small Space,'' \emph{Theory of Computing} 18 (2022), article 1, 1--32. Main derandomization theorem.
\url{https://williamhoza.com/research/papers/CH/}.
{\small\textit{A two-sided hitting-set framework, not an unrestricted logarithmic-seed construction.}}
% END RESEARCH ADDITION SOURCES-163
\end{itemize}

\end{document}
